% ============================================================
\section{Submitted System}
\label{sec:system}
% ============================================================

% Kept short by design. This section describes what we submitted and
% nothing else; everything exploratory lives in Section 6.

Our submission is a three-member majority-vote ensemble over a single
shared prompt. All three members solve the task in the same way --- one
zero-shot call per instance, greedy decoding, a single label token
parsed out of the reply --- and differ only in the underlying model:
one proprietary API model (Claude Opus) and two open-weight models run
locally under \texttt{llama.cpp} (Qwen3.6-35B-A3B and gemma-4-12b-it).
Table~\ref{tab:ensemble} lists the members. Nothing in the system is
trained, fitted, or tuned on labelled data: the only supervision the
system sees is the task's own label definitions and the published
annotation guidelines, which are inserted into the prompt as text.

\begin{table}[t]
\centering
\small
\begin{tabular}{llr}
\toprule
Member & Serving & Open split \\
\midrule
Claude Opus \checkme{exact model id} & API & 0.6640 \\
Qwen3.6-35B-A3B (UD-Q4\_K\_M) & local & 0.6838 \\
gemma-4-12b-it (Q8\_0)        & local & 0.6443 \\
\midrule
Majority vote (submitted)     & ---   & 0.6957 \\
\bottomrule
\end{tabular}
\caption{Ensemble members and their zero-shot accuracy on the open
split ($n{=}253$). Every member uses the identical prompt of
\S\ref{sec:prompt} and the decoding settings of \S\ref{sec:decoding}.}
\label{tab:ensemble}
\end{table}

\subsection{Ensemble and Tie-Breaking}
\label{sec:ensemble}

The three members vote on the five-way label directly; we do not
average scores or probabilities, since one member is served through an
API that does not expose them. With three voters and five classes, a
majority exists whenever any two members agree, and the remaining case
--- three distinct labels --- is resolved in favour of the member with
the highest open-split accuracy (Qwen3.6-35B-A3B). On the open split
this tie-break fires on \checkme{8}/253 instances, and the choice of
tie-breaking member moves accuracy by at most \checkme{1.6}\,pp
(0.6798 with Claude Opus as tie-breaker, 0.6957 with either of the
other two), so the rule is not load-bearing.

We report the ensemble because it is the most accurate configuration we
observed on the open split, but we are explicit about how small the
effect is and how it is obtained. The three members are unanimous on
171/253 open-split instances (67.6\%), and pairwise agreement ranges
from 0.731 (Claude Opus vs.\ Qwen3.6) to 0.795. Voting buys
$+1.2$\,pp of accuracy over the best single member and \emph{costs}
macro-F$_1$ relative to Claude Opus alone (0.395 vs.\ 0.474): the vote
resolves disagreements towards the majority classes, which is precisely
what accuracy rewards and macro-F$_1$ does not. Section~\ref{sec:probe}
returns to this trade-off.
\needrun{Ensemble numbers here were computed offline by majority vote
over the saved per-model prediction files; the voting script is not yet
committed under \texttt{bench/}. Regenerate them from a committed
\texttt{bench/ensemble.py} before submission so the number in the table
and the number in the repository provably match.}

\subsection{Prompt Construction}
\label{sec:prompt}

Each instance in the released data arrives as a fully formed prompt:
role framing, the five label definitions, the Japanese question and
company response, and an instruction to emit the label and nothing
else. We use that prompt unmodified as the carrier and make exactly one
addition: the annotation guidelines are spliced in immediately before
the question--response block, so that the model sees the same
per-label linguistic signals and worked examples --- sentence-final
forms, modality and benefactive markers, conditional expressions ---
that the human annotators were given. The full prompt is reproduced in
Appendix~\ref{app:prompts}. This is a straightforward reading of the
task documentation and we expect it to be common across participating
systems; \S\ref{sec:ablation} quantifies what it contributes and how
that number should be read.

Two aspects of this are worth stating plainly rather than leaving to be
discovered. First, the guidelines are reproduced close to verbatim but
not entirely so: after inspecting errors on the open split we added one
example to the \texttt{+2} block and one signal and one example to the
\texttt{+1} block \checkme{confirm this is the complete list against
the diff of commit \texttt{d517b77}}. These additions are small, and
they were derived from open-split errors only --- no test instance
influenced them --- but the prompt is therefore partly tuned rather
than purely given, and the ablation in \S\ref{sec:ablation} measures
the tuned version. Second, the prompt is English-language instructions
around Japanese content, exactly as released; we did not translate the
instructions into Japanese, and we did not test whether doing so would
help.

We submit a zero-shot system. Few-shot prompting with one exemplar per
label was evaluated on the open split and was worse than zero-shot for
every model we tried once the guidelines were present
(\S\ref{sec:fewshot}), so it does not appear in the submission.

\subsection{Decoding}
\label{sec:decoding}

All members decode greedily (temperature $0$) with a generation budget
of 32 tokens, which is ample for a system asked to emit one of five
short strings. The two local models are served by \texttt{llama.cpp}
with an 8192-token context, all layers offloaded to the GPU, and
extended thinking disabled (\texttt{-{}-reasoning off}); without that
flag both models spend the entire budget on an unfinished reasoning
preamble and emit no label at all. Models that cannot be put into a
non-thinking mode were excluded from the submission for this reason
among others (\S\ref{sec:modelsel}).

The label is recovered from the reply by matching
\texttt{+2\,|\,+1\,|\,0\,|\,-1\,|\,-2} and taking the \emph{last}
match, which favours a trailing final answer over an incidental
numeral earlier in the text. A reply with no match is recorded as
unparsed and counted as an error; we do not retry, resample, or fall
back to a default label, so the reported accuracy is not inflated by
silently repairing failures. In practice this never triggered for the
submitted system: all three members produced 0 unparsed replies on both
the open split (0/253 each) and the test split (0/50 each). The
unparsed rates that do appear in this paper belong to models we did not
submit and are reported with them in Table~\ref{tab:models}.

\subsection{Model Selection}
\label{sec:modelsel}

We evaluated six models zero-shot on the open split under an identical
prompt and identical decoding (Table~\ref{tab:models}) and selected the
submission from those results alone; no test label, and no test
instance, entered the decision. The two open-weight members are the top
two by open-split accuracy (0.6838 and 0.6443) among models that fit in
48\,GB of unified memory and can be made to answer without a reasoning
preamble. Claude Opus was added as a third, architecturally and
operationally independent member: it is the strongest member by
macro-F$_1$ (0.474 vs.\ 0.407 and 0.363) and it is the one whose errors
are least correlated with the others', which is the property a vote
needs. The models we rejected were rejected for stated reasons ---
lower accuracy under the same prompt (shisa-v2-qwen2.5-32b,
Qwen3-30B-A3B-Instruct-2507, gpt-oss-20b) or no non-thinking mode
together with the lowest accuracy in the comparison and a nonzero
unparsed rate (DeepSeek-R1-Distill-Qwen-32B, 7/253 unparsed).

\needrun{Selection was made on open-split point estimates, several of
which are separated by roughly one point on $n{=}253$. The bootstrap
intervals promised in \S\ref{sec:setup} will very likely overlap for
the top three; when they are computed, this subsection should say so
rather than leave the ranking looking sharper than the evidence
supports.}

% ============================================================
\section{Experimental Setup}
\label{sec:setup}
% ============================================================

\paragraph{Data and splits.}
The open split contains 253 question--response pairs with gold labels
and is the only labelled data we used for any decision affecting the
submission. Its label distribution is heavily skewed:
\texttt{+1} 139 (54.9\%), \texttt{0} 72 (28.5\%), \texttt{+2} 29
(11.5\%), \texttt{-1} 11 (4.3\%), \texttt{-2} 2 (0.8\%). Two of the
five classes therefore have single-digit support, which is the reason
we report per-class F$_1$ alongside the aggregate and treat any
claim about \texttt{-1} or \texttt{-2} as indicative only. The test
split contains 50 unlabelled instances and was distributed without gold
labels; it is visibly more balanced than the open split
(\S\ref{sec:annot}) and its instances are markedly shorter, so test and
open-split accuracies are not directly comparable quantities.

No test label was used for prompt design, model selection, tie-breaking,
or any other decision affecting the submission --- and none could have
been, since no test labels were released to participants. The labels we
do have for the test split are our own, produced after submission for
the annotation study of \S\ref{sec:annot}; the protocol below records
where they enter and where they do not.

\paragraph{Models.}
Local models are GGUF quantisations served by \texttt{llama.cpp}
(build \texttt{b10273}, native \texttt{arm64}, Metal backend). Full
identifiers:

\begin{itemize}\itemsep2pt
\item \texttt{unsloth/Qwen3.6-35B-A3B-GGUF},
      \texttt{Qwen3.6-35B-A3B-UD-Q4\_K\_M.gguf} ($\approx$22\,GB; MoE,
      $\approx$3\,B active)
\item \texttt{unsloth/gemma-4-12b-it-GGUF},
      \texttt{gemma-4-12b-it-Q8\_0.gguf} ($\approx$13\,GB; dense 12\,B)
\item \texttt{mradermacher/shisa-v2-qwen2.5-32b-GGUF},
      \texttt{shisa-v2-qwen2.5-32b.Q4\_K\_M.gguf}
\item \texttt{unsloth/Qwen3-30B-A3B-Instruct-2507-GGUF},
      \texttt{Qwen3-30B-A3B-Instruct-2507-Q4\_K\_M.gguf}
\item \texttt{unsloth/gpt-oss-20b-GGUF},
      \texttt{gpt-oss-20b-F16.gguf} (native MXFP4)
\item \texttt{unsloth/DeepSeek-R1-Distill-Qwen-32B-GGUF},
      \texttt{DeepSeek-R1-Distill-Qwen-32B-Q4\_K\_M.gguf}
\end{itemize}

\noindent
The proprietary member is Claude Opus, accessed through the Anthropic
API. \todo{Record the exact model id and snapshot date; the saved
prediction files carry no model metadata, so this must be recovered
from the request logs rather than reconstructed from memory.}
Quantisation is a confound we did not control for: the open-weight
members are compared at different precisions (Q4\_K\_M vs.\ Q8\_0)
because those are the quantisations that fit alongside a 48\,GB budget,
so a gap between two open-weight rows conflates model and precision.

\paragraph{Decoding and reproducibility.}
Temperature $0$ throughout, 32 generation tokens, 8192-token context,
all layers GPU-offloaded, one instance per request with no batching and
no shared state between instances. Greedy decoding makes the sampler
seed irrelevant for the headline runs; where we deliberately sample
(the prompt-robustness probe of \S\ref{sec:sensitivity}) we set
temperature $0.7$ and report results across seeds $\{0,1,2\}$
explicitly. Local runs are deterministic in practice: repeated
executions of the same model and prompt reproduced identical
predictions.

Hardware is a single Apple Silicon machine with 48\,GB of unified
memory \checkme{exact chip}; only one model is resident at a time.
Inference cost for the submitted open-weight members is small: a full
253-instance open-split pass takes 3.1 minutes for Qwen3.6-35B-A3B
(0.74\,s/instance) and 7.3 minutes for gemma-4-12b-it
(1.73\,s/instance), and the 50-instance test pass takes about a minute
each. For contrast, the dense 32\,B reasoning model we did not submit
required 2\,h\,53\,m for the same 253 instances (41\,s/instance), with
7 instances exceeding the 120\,s client timeout. The whole submitted
system can therefore be reproduced on consumer hardware in minutes, API
member aside.

\paragraph{Metrics.}
The official metric is accuracy. We additionally report macro-F$_1$ as
a diagnostic: on a split where the majority class accounts for 54.9\%
of instances, accuracy alone does not distinguish a system that has
learnt the task from one that has collapsed onto the majority class, as
Table~\ref{tab:main} makes clear. The gap is not hypothetical --- an
unweighted fine-tuned encoder baseline reaches 0.5534 accuracy, within
0.4\,pp of always predicting \texttt{+1}, at a macro-F$_1$ of 0.1526.
We report no ordinal metrics, in keeping with the task's categorical
treatment of the label set, but we note that the label set is plainly
ordered and that our own error analysis (\S\ref{sec:probe}) is
concerned almost entirely with adjacent-label confusions.

\paragraph{Significance.}
% Cheap to add, disproportionately increases credibility -- several of
% our deltas are around one point.
We report bootstrap confidence intervals ($B{=}10{,}000$) for accuracy
and use McNemar's test for paired system comparisons. Both are
computed from the saved per-instance predictions, so every comparison
in the paper is paired on the same instances.
\needrun{Bootstrap CIs and McNemar tests are not yet computed. The
comparisons that most need them: ensemble vs.\ best single member
($+1.2$\,pp on $n{=}253$), the top three models against each other, and
every test-split comparison, where $n{=}50$ makes a two-instance
difference worth 4\,pp.}

\paragraph{Evaluation protocol.}
We distinguish throughout between results obtained under official
submission conditions and \emph{post-hoc} analyses. The latter are of
two kinds, and we keep them apart: analyses against the released
open-split gold, and analyses against labels we produced ourselves for
the test split after submission. Table~\ref{tab:protocol} states which
data were available at each stage.

\begin{table}[t]
\centering
\small
\begin{tabular}{ll}
\toprule
Stage & Data used \\
\midrule
Prompt design       & Guidelines + open-split errors \\
Model selection     & Open split gold \\
Official submission & Test split, unlabelled \\
Post-hoc probing (\S\ref{sec:probe}) & Open split gold \\
Annotation study (\S\ref{sec:annot}) & Test split, our labels \\
\bottomrule
\end{tabular}
\caption{Evaluation protocol. No test labels were accessed prior to
submission; none were released to participants. Test-split scores in
\S\ref{sec:annot} are agreement with our own post-hoc annotation, not
official scores, and are labelled as such wherever they appear.}
\label{tab:protocol}
\end{table}

The first row deserves the emphasis. Prompt design was not conditioned
on the guidelines alone: three cues and examples were added after
looking at open-split errors (\S\ref{sec:prompt}). That is within the
rules --- the open split is labelled development data --- but it means
the open-split numbers in Table~\ref{tab:main} are mildly optimistic as
estimates of held-out performance, and we prefer to say so here rather
than have a reader infer it from the commit history.

\todo{Confirm before camera-ready whether official test gold has since
been released. If it has, the test-split results should be recomputed
against it and Table~\ref{tab:protocol} gains a row; if it has not, the
\S\ref{sec:annot} framing as an agreement study rather than an
evaluation must be stated in the abstract as well, not only here.}
